\label{chap.5}
\section{Classifier}
We use DeBERTa-v3-small~\cite{he2021}, fine-tuned to classify each event as malicious or safe with a class-weighted cross-entropy loss. Table~\ref{tab:train} lists the training settings. We keep the epoch with the best validation F1 among those with at most 5\% false positives, and fit a temperature on the validation split so that the output can be used as a calibrated trust score.
\begin{table}[H]
\centering\small
\begin{tabular}{ll}
\toprule
Parameter & Value\\
\midrule
Base model & microsoft/deberta-v3-small\\
Maximum length & 512 tokens\\
Batch size & 4, gradient accumulation 4\\
Learning rate, epochs, warm-up & $2\times10^{-5}$, 2, 0.1\\
Precision & fp16 with fp32 master weights\\
Seeds & 42, 2, 3, 4, 5 (five runs reported)\\
Hardware & RTX 3050 Laptop GPU\\
Training time & about 4.3 h (stateful), 1.3 h (stateless)\\
\bottomrule
\end{tabular}
\caption{Training configuration.}
\label{tab:train}
\end{table}

\section{Training Data}
To build the training data we replay every training session through \texttt{session.py}, so each event text carries the same session signals that MAPIS will compute at run time. Only the event with the injection and the final action are positives; the surrounding events are labelled safe. We also add hard negatives, clean twins and clean windows next to the attacks. The training and validation parts of MAPIS-Adaptive are included as adversarial training data, about 52,000 chunks, 34\% more than before.

\section{Stateless Ablation}
To see what the session memory adds, we trained the same model on the same data but gave it only the current event (\texttt{--no-context}). This is the \emph{stateless DeBERTa} used all through Chapter~6, and it is the fair ``same model, no memory'' comparison.

\section{Implementation}
\begin{table}[H]
\centering\small
\begin{tabular}{p{4.4cm}p{8.3cm}}
\toprule
Component & Implementation\\
\midrule
Shield & \texttt{backend/core/shield.py}: normalise, trust anchor, suspended channels, score, fuse, tier\\
Session state & \texttt{session.py}, \texttt{features.py}, \texttt{context.py}, \texttt{store.py} (Redis or memory)\\
Detector & \texttt{detector.py}: transformer detector with chunking, regex fallback\\
API & FastAPI: \texttt{/inspect}, \texttt{/pipeline/run}, events, quarantine review, WebSocket, replay, results\\
Testbed & \texttt{testbed.py}, \texttt{graph.py} (LangGraph), \texttt{autogen\_testbed.py} (AutoGen)\\
Evaluation & \texttt{benchmark.py}, \texttt{evaluate.py}, \texttt{stateless\_sweep.py}, \texttt{live\_agentdojo.py}, \texttt{baselines.py}\\
Interface & React operator console in \texttt{frontend/}\\
Tests & 33 pytest tests: shield, API, train/serve parity, AutoGen tap, multi-hop attacks against twins, tampering checks\\
\bottomrule
\end{tabular}
\caption{Implementation overview.}
\end{table}

\section{Development History and Fixes}
\label{sec:history}
We did not get to the final system in one go. Each of the failures below was measured, and each one changed the design.
\begin{enumerate}
\item \textbf{Placeholder numbers.} Early slides quoted results that were placeholders, not measurements. All numbers were rebuilt and are now taken from the logs.
\item \textbf{Truncation.} The scored event was placed last in the text and cut at 384 tokens, hiding 83\% of AgentDojo injections. The event now comes first, long events are chunked, and trust is the minimum over chunks.
\item \textbf{The model ignored the session.} An early model gained only 2.3 points from state. Moving the session signals into the classifier text, built by the same code in training and at run time, fixed this.
\item \textbf{Misses on override and physical-safety attacks.} Nothing travels across hops there, and the model had memorised tool names. We added the unrequested-action and claim-echo signals, an escalation term and more action families.
\item \textbf{Multi-hop false positives of 10 to 17\%.} Two causes: contradictory labels in our generator (the same text labelled malicious in attacks and benign in twins) and memorised tool names. We changed the label policy, added eight further override and physical action kinds as training families, and drew carrier text from a larger pool. The multi-hop false-positive rate fell to about 2\%.
\item \textbf{Evasive attacker.} Against the adaptive set the first model reached 64.6\% recall. Canonicalisation, fragment matching and adversarial training raised it to 100\% on all six techniques.
\item \textbf{A stale old-code run.} One run silently used old code because a pull failed. The pipeline now aborts at its first step if the code is out of date.
\item \textbf{Plain values on the team-written set.} MAPIS caught only 35 to 45\% of the team's attacks: provenance did not follow changed values, privilege flags or destinations relayed by an agent. The tampering checks of Section~\ref{sec:tamper} raised this to 95\% without retraining (Section~\ref{sec:harder}).
\end{enumerate}
